
\documentclass[11pt,letterpaper]{article}
\usepackage[utf8]{inputenc}
\usepackage[T1]{fontenc}
\usepackage[spanish]{babel}
\usepackage{lmodern,amsmath,amssymb,geometry,xcolor,tikz,tcolorbox,fancyhdr,enumitem,lastpage,array}
\usetikzlibrary{arrows.meta,calc,positioning}
\tcbuselibrary{skins,breakable}
\geometry{margin=1.85cm,top=2.3cm,bottom=2cm}
\definecolor{azul}{RGB}{20,65,115}\definecolor{verde}{RGB}{0,130,90}\definecolor{morado}{RGB}{116,53,166}\definecolor{naranja}{RGB}{225,103,33}
\setlength{\parindent}{0pt}\setlength{\parskip}{5pt}\setlength{\headheight}{15pt}
\pagestyle{fancy}\fancyhf{}\lhead{\textcolor{azul}{\bfseries Física -- Educación Diversificada}}\rhead{\textcolor{azul}{J. Arias}}\cfoot{\thepage\ de \pageref{LastPage}}\renewcommand{\headrulewidth}{0.5pt}
\newtcolorbox{bloque}[1]{enhanced,breakable,colback=blue!3,colframe=azul,colbacktitle=azul,coltitle=white,title=\bfseries #1,arc=2mm,boxrule=.65pt}
\newtcolorbox{problema}[1]{enhanced,breakable,colback=white,colframe=azul!50,colbacktitle=blue!9,coltitle=azul,title=\bfseries Ejercicio #1,arc=2mm,boxrule=.6pt}
\newcommand{\espacio}[1]{\par\vspace{#1}\par}
\newcommand{\renglon}{\par\textcolor{morado}{\bfseries Respuesta:}\ \dotfill\par}
\newcommand{\vehiculo}[5]{% x y color nombre velocidad
\draw[fill=#3!12,draw=#3,thick,rounded corners=2pt] (#1,#2) rectangle ++(1.3,.38);
\fill[black!70] ($(#1,#2)+(.25,-.03)$) circle (.085);
\fill[black!70] ($(#1,#2)+(1.05,-.03)$) circle (.085);
\node[font=\bfseries\scriptsize] at ($(#1,#2)+(.65,.19)$) {#4};
\node[above,font=\scriptsize] at ($(#1,#2)+(.65,.40)$) {#5};}
\newcommand{\flechaD}[3]{\draw[-{Stealth[length=2.5mm]},very thick,#3] (#1,#2)--++(1.2,0);}
\newcommand{\flechaI}[3]{\draw[-{Stealth[length=2.5mm]},very thick,#3] (#1,#2)--++(-1.2,0);}
\newcommand{\cuatroautos}[8]{%
\begin{center}\begin{tikzpicture}[scale=.9]
\foreach \y in {0,1.05,2.1,3.15}{\draw[gray!50] (.2,\y)--(10.4,\y);}
\vehiculo{1}{3.25}{blue}{A}{#1}\vehiculo{6.2}{2.2}{orange}{B}{#3}
\vehiculo{1}{1.15}{green!60!black}{C}{#5}\vehiculo{6.2}{.1}{purple}{D}{#7}
\csname flecha#2\endcsname{2.45}{3.43}{blue}
\csname flecha#4\endcsname{6.05}{2.38}{orange}
\csname flecha#6\endcsname{2.45}{1.33}{green!60!black}
\csname flecha#8\endcsname{6.05}{.28}{purple}
\end{tikzpicture}\end{center}}
\begin{document}
\begin{center}{\LARGE\bfseries\color{azul} PRÁCTICA DE FÍSICA}\\[3pt]
{\Large Movimiento relativo en una dimensión}\\[4pt]
{\color{verde}\bfseries Vehículos, pasajeros y observadores}\end{center}
\begin{tcolorbox}[colback=blue!3,colframe=azul,arc=2mm]
\textbf{Estudiante:}\ \hrulefill\\[8pt]
\textbf{Sección:}\ \underline{\hspace{3cm}}\hfill\textbf{Fecha:}\ \underline{\hspace{3.5cm}}
\end{tcolorbox}
\begin{bloque}{Indicaciones}
Resuelva cada apartado mostrando datos, sustitución y cálculo. Use \textbf{positivo hacia el este o la derecha} y \textbf{negativo hacia el oeste o la izquierda}. Exprese la velocidad relativa con signo, unidad y sentido. No se incluyen soluciones.
\end{bloque}
\section*{I. Interpretación de observadores}
\begin{problema}{1}
Un autobús viaja hacia el este a $60\,\mathrm{km/h}$. Una pasajera está sentada y una persona permanece inmóvil en la acera.
\begin{enumerate}[label=\alph*)]
\item Indique la velocidad de la pasajera respecto al autobús.
\item Indique la velocidad de la pasajera respecto al suelo.
\item Determine la velocidad de la persona de la acera respecto al autobús.
\item Explique quién observa a la pasajera en reposo.
\end{enumerate}
\espacio{3cm}\renglon
\end{problema}
\newpage
\begin{problema}{2}
Un automóvil A se mueve al este a $45\,\mathrm{km/h}$ y otro B al este a $75\,\mathrm{km/h}$.
\begin{enumerate}[label=\alph*)]\item Calcule $v_{A/B}$.\item Calcule $v_{B/A}$.\item Indique la rapidez relativa entre ambos.\item Interprete los signos.\end{enumerate}
\espacio{3cm}\renglon
\end{problema}
\clearpage
\section*{II. Vehículos en movimiento}
\begin{problema}{3}
Observe los cuatro vehículos y determine las velocidades relativas solicitadas.
\cuatroautos{$40\,\mathrm{km/h}$}{D}{$50\,\mathrm{km/h}$}{I}{$30\,\mathrm{km/h}$}{D}{$35\,\mathrm{km/h}$}{D}
\begin{enumerate}[label=\alph*)]
\item $v_{A/B}$ \qquad \item $v_{A/C}$ \qquad \item $v_{D/B}$
\item $v_{B/C}$ \qquad \item $v_{C/D}$ \qquad \item $v_{B/A}$
\end{enumerate}
\espacio{4cm}\renglon
\end{problema}
\newpage
\begin{problema}{4}
Para los vehículos del siguiente esquema, calcule cada velocidad relativa e indique su sentido.
\cuatroautos{$70\,\mathrm{km/h}$}{D}{$45\,\mathrm{km/h}$}{D}{$35\,\mathrm{km/h}$}{I}{$60\,\mathrm{km/h}$}{I}
\begin{enumerate}[label=\alph*)]
\item $v_{A/B}$ \qquad \item $v_{C/D}$ \qquad \item $v_{D/A}$
\item $v_{B/C}$ \qquad \item $v_{A/D}$ \qquad \item $v_{C/B}$
\end{enumerate}
\espacio{4cm}\renglon
\end{problema}
\clearpage
\begin{problema}{5}
Considere el siguiente conjunto de vehículos.
\cuatroautos{$25\,\mathrm{km/h}$}{I}{$55\,\mathrm{km/h}$}{I}{$80\,\mathrm{km/h}$}{D}{$20\,\mathrm{km/h}$}{D}
\begin{enumerate}[label=\alph*)]
\item $v_{A/C}$ \qquad \item $v_{B/A}$ \qquad \item $v_{C/D}$
\item $v_{D/B}$ \qquad \item $v_{A/D}$ \qquad \item $v_{C/B}$
\end{enumerate}
\espacio{4cm}\renglon
\end{problema}
\newpage
\begin{problema}{6}
Dos motocicletas se aproximan una a la otra. La motocicleta A se mueve hacia el este a $54\,\mathrm{km/h}$ y la motocicleta B hacia el oeste a $72\,\mathrm{km/h}$.
\begin{enumerate}[label=\alph*)]
\item Determine $v_{A/B}$ y $v_{B/A}$.
\item Determine la rapidez relativa.
\item Si la separación inicial es $630\,\mathrm{m}$ y mantienen velocidades constantes, determine el tiempo hasta encontrarse.
\end{enumerate}
\espacio{4.5cm}\renglon
\end{problema}
\clearpage
\section*{III. Personas dentro y fuera de un autobús}
\begin{problema}{7}
Un autobús avanza al este a $72\,\mathrm{km/h}$. Marta camina hacia adelante a $4\,\mathrm{km/h}$ respecto al autobús; Pedro camina hacia atrás a $3\,\mathrm{km/h}$ respecto al autobús; Sara permanece sentada.
\begin{center}\begin{tikzpicture}[scale=.85]
\draw[rounded corners=5pt,fill=blue!8,draw=azul,very thick] (0,0) rectangle (10,2.15);
\foreach \x in {1,3,5,7}{\draw[fill=white] (\x,1.25) rectangle ++(1.3,.65);}
\foreach \x in {1.6,8.4}{\fill[black!75] (\x,-.07) circle (.32);}
\node at (5,.5) {\bfseries AUTOBÚS};
\draw[-{Stealth},very thick,verde] (10.2,1)--(11.7,1);
\node[above] at (5,2.3) {$72\,\mathrm{km/h}$ hacia el este};
\node[font=\scriptsize] at (1.5,1.05) {Marta $+4$};
\node[font=\scriptsize] at (4,1.05) {Pedro $-3$};
\node[font=\scriptsize] at (7.8,1.05) {Sara $0$};
\end{tikzpicture}\end{center}
\begin{enumerate}[label=\alph*)]
\item Calcule las velocidades de Marta, Pedro y Sara respecto al suelo.
\item Determine $v_{M/P}$ y $v_{P/M}$.
\item ¿Qué velocidad observa Sara para Marta?
\item ¿Qué velocidad observa Pedro para Sara?
\item ¿Qué velocidad observa Sara para alguien inmóvil en la acera?
\end{enumerate}
\espacio{5cm}\renglon
\end{problema}
\clearpage
\begin{problema}{8}
Un autobús se mueve hacia el oeste a $50\,\mathrm{km/h}$. Laura camina dentro del autobús hacia el este a $6\,\mathrm{km/h}$ respecto al vehículo. Un ciclista se desplaza por la carretera hacia el oeste a $15\,\mathrm{km/h}$.
\begin{enumerate}[label=\alph*)]
\item Determine la velocidad de Laura respecto al suelo.
\item Determine la velocidad del ciclista respecto al autobús.
\item Determine la velocidad de Laura respecto al ciclista.
\item Determine la velocidad del autobús respecto a Laura.
\end{enumerate}
\espacio{5cm}\renglon
\end{problema}
\begin{problema}{9}
Un autobús viaja hacia el este a $65\,\mathrm{km/h}$. Una persona camina hacia la parte trasera del autobús a $5\,\mathrm{km/h}$ respecto a este. En la acera, otra persona camina hacia el este a $4\,\mathrm{km/h}$.
\begin{enumerate}[label=\alph*)]
\item Determine la velocidad de la persona dentro del autobús respecto al suelo.
\item Calcule la velocidad de la persona dentro del autobús respecto a quien camina en la acera.
\item Determine la velocidad de quien camina en la acera respecto al autobús.
\end{enumerate}
\espacio{4.5cm}\renglon
\end{problema}
\clearpage
\section*{IV. Trenes, bandas y escaleras mecánicas}
\begin{problema}{10}
Un tren se desplaza hacia el este a $24\,\mathrm{m/s}$. Una pasajera camina hacia atrás a $2\,\mathrm{m/s}$ respecto al tren. Un corredor en el andén se mueve al este a $5\,\mathrm{m/s}$.
\begin{enumerate}[label=\alph*)]
\item Calcule la velocidad de la pasajera respecto al suelo.
\item Determine la velocidad de la pasajera respecto al corredor.
\item Determine la velocidad del corredor respecto al tren.
\item Calcule la velocidad del tren respecto a la pasajera.
\end{enumerate}
\espacio{4.5cm}\renglon
\end{problema}
\begin{problema}{11}
Una banda transportadora se mueve hacia la derecha a $1,5\,\mathrm{m/s}$. Daniel camina sobre ella hacia la derecha a $1,2\,\mathrm{m/s}$ respecto a la banda. Andrea camina sobre la misma banda hacia la izquierda a $0,8\,\mathrm{m/s}$ respecto a la banda.
\begin{enumerate}[label=\alph*)]
\item Calcule la velocidad de Daniel respecto al suelo.
\item Calcule la velocidad de Andrea respecto al suelo.
\item Determine la velocidad de Daniel respecto a Andrea.
\item Determine la velocidad de Andrea respecto a Daniel.
\end{enumerate}
\espacio{4.5cm}\renglon
\end{problema}
\clearpage
\begin{problema}{12}
Una escalera mecánica se mueve hacia la derecha a $0,8\,\mathrm{m/s}$. Una persona camina hacia la izquierda a $1,1\,\mathrm{m/s}$ respecto a la escalera. Otra persona permanece de pie sobre ella.
\begin{enumerate}[label=\alph*)]
\item Determine la velocidad de cada persona respecto al suelo.
\item Calcule la velocidad de la primera persona respecto a la segunda.
\item Indique si la primera persona avanza a la derecha o a la izquierda respecto al suelo.
\end{enumerate}
\espacio{5cm}\renglon
\end{problema}
\section*{V. Desafíos de aplicación}
\begin{problema}{13}
Un autobús viaja al este a $54\,\mathrm{km/h}$. Un estudiante camina dentro del autobús hacia atrás a $1,5\,\mathrm{m/s}$ respecto al vehículo. Una persona está inmóvil en la parada.
\begin{enumerate}[label=\alph*)]
\item Convierta la velocidad del autobús a $\mathrm{m/s}$.
\item Determine la velocidad del estudiante respecto al suelo en $\mathrm{m/s}$.
\item Determine la velocidad de la persona de la parada respecto al estudiante.
\item Exprese la rapidez del estudiante respecto al suelo en $\mathrm{km/h}$.
\end{enumerate}
\espacio{4cm}\renglon
\end{problema}
\clearpage
\begin{problema}{14}
Un tren se desplaza al oeste a $18\,\mathrm{m/s}$. Una persona camina hacia el este dentro del tren a $3\,\mathrm{m/s}$ respecto al tren. Un automóvil avanza hacia el oeste a $12\,\mathrm{m/s}$ respecto al suelo.
\begin{enumerate}[label=\alph*)]
\item Calcule la velocidad de la persona respecto al suelo.
\item Calcule la velocidad de la persona respecto al automóvil.
\item Calcule la velocidad del automóvil respecto al tren.
\item ¿Con qué velocidad debería caminar la persona respecto al tren para quedar en reposo respecto al suelo?
\end{enumerate}
\espacio{5cm}\renglon
\end{problema}
\begin{tcolorbox}[colback=orange!5,colframe=naranja,title=\bfseries Revisión final,arc=2mm]
$\square$ Asigné correctamente los signos.\qquad $\square$ Identifiqué el observador.\\
$\square$ Conservé las unidades.\qquad $\square$ Interpreté el sentido de cada resultado.
\end{tcolorbox}
\end{document}